\chapter{QNX-Compatible C Library}

\section{Purpose}

QSOE provides a C library layer that exposes QNX Neutrino-compatible
IPC and threading primitives.  This library is linked by all user-space
programs (including taskman itself) and translates QNX API calls into
the appropriate seL4 system calls and capability operations.

It is built on top of musl libc, which provides standard POSIX/C
functions, and extends it with the QNX-specific interfaces.

\section{Core IPC functions}

\subsection{Channels}

\begin{lstlisting}[language=C]
int ChannelCreate(unsigned flags);
int ChannelDestroy(int chid);
\end{lstlisting}

A \emph{channel} is the server-side IPC endpoint.  Creating a channel
allocates an seL4 Endpoint capability and registers it in the process's
channel table.

\subsection{Connections}

\begin{lstlisting}[language=C]
int ConnectAttach(uint32_t nd, pid_t pid, int chid,
                  unsigned index, int flags);
int ConnectDetach(int coid);
\end{lstlisting}

A \emph{connection} is the client-side handle.  \texttt{ConnectAttach}
obtains a badged copy of the target channel's Endpoint capability,
enabling the client to send messages.

\subsection{Messages}

\begin{lstlisting}[language=C]
int MsgSend(int coid, const void *smsg, int sbytes,
            void *rmsg, int rbytes);
int MsgReceive(int chid, void *msg, int bytes,
               struct _msg_info *info);
int MsgReply(int rcvid, int status,
             const void *msg, int bytes);
int MsgSendPulse(int coid, int priority,
                 int code, int value);
\end{lstlisting}

These form the fundamental Send/Receive/Reply triad.  The synchronous
variants map directly to \texttt{seL4\_Call}, \texttt{seL4\_Recv}, and
\texttt{seL4\_Reply}.  Pulses use seL4 Notification objects.

\section{Threading}

\begin{lstlisting}[language=C]
int ThreadCreate(pid_t pid, void *(*func)(void *),
                 void *arg, const struct _thread_attr *attr);
int ThreadDestroy(int tid, int priority, void *status);
\end{lstlisting}

Thread creation allocates a new seL4 TCB (Thread Control Block),
configures its scheduling parameters, and starts execution at the
specified function.

\section{Relationship to musl libc}

Standard C/POSIX functions (\texttt{printf}, \texttt{malloc},
\texttt{open}, \texttt{read}, etc.) are provided by musl libc, patched
to route system calls through seL4 IPC to \texttt{taskman} or the
appropriate resource manager rather than through Linux \texttt{ecall}
traps.
